\section{Método de pesquisa}
\label{sec:metodo}

A pesquisa foi estruturada como revisão integrativa sistematizada, com apoio bibliométrico e síntese temática, reunindo estudos de diferentes delineamentos sob busca, seleção e síntese explicitados \parencite{whittemore_revisaointegrativa_2005,snyder_revisaometodologia_2019}. O apoio bibliométrico organizou o corpus e orientou a síntese temática \parencite{donthu_bibliometria_2021}, com transparência e reprodutibilidade destacadas por \textcite{hu_revisao_sintese_2026}. O escopo declarado distingue esses procedimentos das etapas de pré-registro, dupla revisão, avaliação de risco de viés, elegibilidade integral em texto completo e metanálise, ausentes deste desenho.

\subsection{Conjuntos analíticos}

A análise foi organizada em dois níveis complementares. A camada bibliométrica ampliada utilizou os 372 registros do estrato bibliométrico forte da busca principal, empregados para analisar fontes, coocorrência de palavras-chave, mapa temático e evolução entre 2010--2020 e 2021--2026 (completude dos campos na Seção~\ref{sec:panorama}). A síntese temática utilizou inicialmente 104 registros e passou a 121 após a incorporação de 17 registros identificados pela busca complementar de sensibilidade para inteligência artificial e aprendizado de máquina. A comparação por identificador único mostrou que 109 dos 121 registros do núcleo temático vigente também pertencem ao estrato bibliométrico de 372, 12 são exclusivos do núcleo temático vigente e 263 são exclusivos da camada bibliométrica; os conjuntos permanecem parcialmente sobrepostos e não foram fundidos, para não alterar artificialmente a estrutura temática da busca principal.

A Tabela~\ref{tab:glossariocamadas} padroniza os rótulos usados no artigo e evita que diferentes etapas sejam tratadas como um único conjunto.

\begin{table}[htbp]
\centering
\caption{Glossário dos conjuntos analíticos da revisão}
\label{tab:glossariocamadas}
\small
\begin{tabularx}{\textwidth}{>{\RaggedRight\arraybackslash}p{3.3cm} >{\raggedleft\arraybackslash}p{1.6cm} Y}
\toprule
Rótulo adotado & Registros & Definição \\
\midrule
Núcleo analítico revisado & 3.678 & Registros retidos após pré-triagem e resolução dos casos de dúvida. \\
Estrato bibliométrico forte & 372 & Subconjunto da busca principal usado exclusivamente para a estrutura bibliométrica do campo. \\
Candidatos à auditoria qualitativa & 137 & Registros centrais após o alinhamento determinístico às perguntas de pesquisa. \\
Núcleo temático original & 104 & Registros mantidos após auditoria qualitativa estruturada de título, resumo, palavras-chave e campos extraídos. \\
Núcleo temático vigente & 121 & Núcleo original acrescido de 17 registros incorporados pela busca de sensibilidade IA/ML. \\
\bottomrule
\end{tabularx}
\fonteautor
\end{table}

A Tabela~\ref{tab:alinhamentorq} explicita a correspondência entre a pergunta central, as questões específicas, os critérios de seleção, os campos extraídos e os resultados apresentados. O contexto público universitário e os métodos de apoio à decisão foram campos de caracterização analítica, e não requisitos obrigatórios de inclusão.

\begin{table}[htbp]
\centering
\caption{Matriz de alinhamento entre perguntas, seleção, extração e resultados}
\label{tab:alinhamentorq}
\scriptsize
\begin{tabularx}{\dimexpr\textwidth-2pt\relax}{>{\RaggedRight\arraybackslash}p{1.8cm} >{\RaggedRight\arraybackslash}p{2.7cm} >{\RaggedRight\arraybackslash}p{2.3cm} >{\RaggedRight\arraybackslash}p{4.1cm} Y}
\toprule
Elemento & Formulação sintética & Critério de seleção & Campo de extração & Resultado correspondente \\
\midrule
RQ0 & Integração entre sustentabilidade, manutenção, gestão predial e apoio à decisão; elementos aplicáveis à priorização pública universitária. & Objeto predial e sustentabilidade obrigatórios; decisão e contexto como caracterização. & \texttt{rq0\_\allowbreak confirmada}; contribuição; critérios; métodos; contexto; dimensões. & Síntese integrada nas Seções~\ref{sec:criterios} a~\ref{sec:matriz}. \\
RQ1 & Dimensões de sustentabilidade identificadas. & Sustentabilidade obrigatória. & \texttt{dimensoes\_\allowbreak sustentabilidade\_\allowbreak identificadas\_\allowbreak leitura}. & Seção~\ref{sec:criterios} e Gráfico~\ref{fig:dimensoes}. \\
RQ2 & Critérios de priorização identificados. & Decisão e priorização como caracterização. & \texttt{criterios\_\allowbreak identificados\_\allowbreak leitura}. & Seção~\ref{sec:criterios} e Tabela~\ref{tab:criteriospriorizacao}. \\
RQ3 & Métodos de apoio à decisão identificados. & Método de decisão como caracterização. & \texttt{metodos\_\allowbreak decisao\_\allowbreak identificados\_\allowbreak leitura}. & Seção~\ref{sec:metodos} e Gráfico~\ref{fig:metodos}. \\
RQ4 & Contextos de edificação identificados. & Contexto como caracterização, sem recorte isolado. & \texttt{contexto\_\allowbreak edificacao\_\allowbreak identificado\_\allowbreak leitura}. & Seção~\ref{sec:aplicabilidade} e Tabela~\ref{tab:contexto}. \\
RQ5 & Lacunas registradas para instituições públicas de ensino superior. & Contexto público/\allowbreak universitário como reforço não obrigatório. & \texttt{lacuna\_\allowbreak ou\_\allowbreak limite\_\allowbreak identificado\_\allowbreak leitura}; contribuição do artigo. & Seções~\ref{sec:aplicabilidade} e \ref{sec:consideracoes}. \\
RQ6 & Efeito da busca de sensibilidade sobre a identificação de IA e aprendizado de máquina. & Técnica computacional confirmada, objeto predial e sustentabilidade; termos ambíguos exigem contexto adicional. & Classe de evidência RQ6; técnica identificada; contexto; decisão de incorporação. & Seções~\ref{sec:buscaia}, \ref{sec:metodos} e \ref{sec:consideracoes}. \\
\bottomrule
\end{tabularx}
\fonteautor
\end{table}

\subsection{Estratégia de busca}

A busca principal abrangeu publicações de 2010 a 2026 nas bases Scopus, Web of Science e Crossref, com 13 consultas executadas em 8 de julho de 2026 e enriquecimento manual da Scopus em 9 de julho, incorporando cinco registros do arquivo manual. Na Scopus, quatro expressões booleanas foram executadas em \texttt{TITLE-ABS-KEY}, sem filtros de idioma ou tipo documental; na Web of Science, quatro buscas manuais equivalentes foram registradas no campo \texttt{TS (Topic)}. No Crossref, cinco consultas por \texttt{query.bibliographic} usaram filtro de data entre 2010-01-01 e 2026-12-31 e limite de 200 registros por consulta; por ordenar resultados por relevância, o Crossref foi tratado como fonte complementar de metadados, e não como equivalente às expressões booleanas das demais bases \parencite{crossref_restapi_2026}.

A busca complementar de sensibilidade foi executada em 12 de julho de 2026 nas mesmas três fontes, sem filtros adicionais além do período de publicação (Tabela~\ref{tab:estrategiabusca}). A soma de 18.846 ocorrências entre as duas rodadas corresponde apenas ao volume operacional antes da deduplicação e não constitui corpus homogêneo, porque a busca de sensibilidade foi deliberadamente enriquecida para IA e aprendizado de máquina.

\begin{table}[htbp]
\centering
\caption{Estratégia de busca consolidada por rodada e base}
\label{tab:estrategiabusca}
\scriptsize
\begin{tabularx}{\textwidth}{>{\RaggedRight\arraybackslash}p{2.1cm} >{\RaggedRight\arraybackslash}p{2.1cm} >{\RaggedRight\arraybackslash}p{1.6cm} >{\RaggedRight\arraybackslash}p{3.0cm} >{\raggedleft\arraybackslash}p{1.3cm} Y}
\toprule
Rodada & Base & Data & Campo e consultas & Registros & Documentação \\
\midrule
Principal & Scopus & 08--09/07/2026 & 4 em \texttt{TITLE-ABS-KEY}; enriquecimento por EID & 9.438 & Strings preservadas; 9.433 da API e cinco do export manual. \\
Principal & Web of Science & 08/07/2026 & 4 em \texttt{TS (Topic)} & 1.680 & Strings preservadas; quarta exportação recontada em 503. \\
Principal & Crossref & 08/07/2026 & 5 em \texttt{query.bibliographic} & 1.000 & Consultas preservadas; 200 resultados por relevância. \\
\cmidrule(lr){1-6}
\multicolumn{4}{r}{Subtotal da busca principal} & 12.118 & 13 consultas, antes da deduplicação. \\
\midrule
Sensibilidade IA/ML & Scopus & 12/07/2026 & 1 em \texttt{TITLE-ABS-KEY} & 3.169 & String integral documentada; sem filtros adicionais além do período. \\
Sensibilidade IA/ML & Web of Science, All Databases & 12/07/2026 & 1 em \texttt{Topic (TS)}, exportada em duas partes & 1.559 & String integral documentada; registros indexados na Core Collection. \\
Sensibilidade IA/ML & Crossref & 12/07/2026 & 10 em \texttt{query.bibliographic} & 2.000 & Consultas preservadas; 1.993 após deduplicação interna. \\
\cmidrule(lr){1-6}
\multicolumn{4}{r}{Subtotal da busca de sensibilidade} & 6.728 & 12 consultas, antes da deduplicação. \\
\bottomrule
\end{tabularx}
\par\smallskip{\footnotesize Nota: a soma operacional das rodadas é 18.846 ocorrências, mas não representa um único corpus bruto. Fonte: elaborado pelo autor a partir dos logs e arquivos exportados.\par}
\end{table}

Na Web of Science, a associação entre cada arquivo exportado e as quatro strings principais foi informada de memória pelo pesquisador, não reconstruída de registro contemporâneo. Não houve busca piloto, uso de estudos-semente ou pré-registro; o protocolo e a matriz conceitual foram elaborados antes da coleta principal, e a RQ6 foi formalizada posteriormente para testar a insuficiência temática identificada. Os critérios de seleção combinaram aderência ao ambiente construído, manutenção ou gestão predial e sustentabilidade, com método de decisão, contexto público universitário e tipologia da edificação como campos de caracterização analítica (Tabela~\ref{tab:criteriosincluexclu}).

\begin{table}[htbp]
\centering
\caption{Critérios de seleção e caracterização do corpus}
\label{tab:criteriosincluexclu}
\footnotesize
\begin{tabularx}{\textwidth}{>{\RaggedRight\arraybackslash}p{3.4cm} >{\RaggedRight\arraybackslash}p{2.7cm} Y}
\toprule
Critério & Aplicação & Descrição operacional \\
\midrule
Objeto predial & Inclusão & Manutenção predial, gestão de facilidades, gestão de ativos construídos ou operação de edifícios. \\
Sustentabilidade & Inclusão & Desempenho ambiental, ciclo de vida, energia, critérios econômicos, sociais, técnicos ou institucionais. \\
Período de publicação & Inclusão & Publicação entre 2010 e 2026. \\
Decisão e priorização & Caracterização & Identificação de MCDM, MCDA, AHP, TOPSIS, ANP, otimização, pontuação e outros instrumentos de apoio à decisão. \\
Contexto da edificação & Caracterização & Identificação de universidades, \textit{campus}, edifícios públicos, escolas, hospitais e outras tipologias. \\
Escopo externo ao ambiente construído & Exclusão & Aplicações industriais, transporte, infraestrutura linear ou materiais sem relação com operação e manutenção de edificações. \\
Duplicidade & Consolidação & Fusão por DOI normalizado ou título normalizado, com preservação das bases e consultas de origem. \\
\bottomrule
\end{tabularx}
\fonteautor
\end{table}

\subsection{Deduplicação e preservação da proveniência}

A deduplicação foi executada após padronizar os metadados das três fontes, com o DOI convertido para minúsculas e sem prefixos de URL, espaços ou ponto final, e o título normalizado quanto a caixa, acentuação e pontuação, excluindo do agrupamento títulos com menos de oito caracteres, preparação adaptada dos problemas de interoperabilidade documentados por \textcite{rathbone_deduplicacao_2015}, sem adoção do algoritmo SRA-DM. Registros com o mesmo DOI normalizado foram agrupados; na ausência de DOI, títulos normalizados idênticos só foram agrupados sem conflito entre identificadores presentes, critério mais restritivo que a abordagem baseada em regras de \textcite{jiang_deduplicacao_regras_2014}. Para cada grupo fundido, preservaram-se bases, consultas de origem e identificadores brutos, mantendo o resumo mais longo disponível e selecionando os demais metadados por completude e prioridade de fonte.

A aplicação dessas regras consolidou 12.118 ocorrências em 9.542 registros únicos, removendo 2.576 repetições. Dos 1.808 grupos com duplicidade, 1.635 foram agrupados por DOI e 173 por título normalizado, com 98 conflitos de identificadores mantidos separados. A Tabela S1 do material suplementar apresenta os indicadores completos.

\subsection{Triagem e auditoria}

A pré-triagem dos 9.542 registros foi uma classificação automática determinística, sem uso de modelo de linguagem, baseada no casamento de expressões em título e resumo, distribuindo os registros em cinco classes (relevante, complementar, dúvida, irrelevante e sem resumo). Uma amostra estratificada de 100 registros (semente 42, 1,0\% do corpus único) foi lida individualmente por um único avaliador, que alterou 34 classificações; os demais 9.442 registros não foram lidos individualmente nessa auditoria amostral. O corte inicial reteve 3.786 registros, dos quais 3.472 relevantes e 314 casos de dúvida com apenas o bloco de objeto predial; esses 314 casos integravam o universo de 4.276 dúvidas reavaliadas individualmente, do qual 206 foram considerados relevantes e 4.070 descartados, resultando em um núcleo analítico revisado de 3.678 registros validado por script em R, sem refazer a triagem. A documentação preserva decisões, regras, níveis de confiança e justificativas, mas não registra revisão independente nem resolução formal de divergências entre avaliadores.

Os 3.678 registros passaram por duas camadas determinísticas em R, sem modelo de linguagem e sem leitura de texto completo. A primeira classificou 372 como estrato bibliométrico forte, 830 descritivos, 157 contextuais e 2.319 descarte; a segunda pontuou o alinhamento às perguntas de pesquisa, classificando 137 como centrais, 651 secundários, 375 descritivos, 157 contextuais e 2.358 excluídos, encaminhando os 137 centrais à auditoria qualitativa. Essa auditoria examinou título, resumo, palavras-chave e campos extraídos, sem leitura integral, mantendo 104 no núcleo temático original; entre os 33 restantes, 21 foram secundários, três apenas para mapeamento e nove excluídos, um deles previamente marcado para verificação pontual em texto completo. Não houve segundo avaliador independente. Os arquivos da auditoria não identificam ferramenta de IA, modelo, versão ou prompt, e o campo de compatibilidade das planilhas não constitui evidência de uso do ASReview. A Tabela~\ref{tab:funilselecao} preserva a rastreabilidade da busca principal até o núcleo original de 104 registros, e a Figura~\ref{fig:funil} integra esse percurso à busca de sensibilidade, explicitando a convergência no núcleo temático vigente de 121 registros.

{\setlength{\tabcolsep}{2pt}
\begin{table}[htbp]
\centering
\caption{Rastreabilidade do funil de seleção da busca principal}
\label{tab:funilselecao}
\tiny
\begin{tabularx}{\textwidth}{>{\RaggedRight\arraybackslash}p{1.7cm} >{\raggedleft\arraybackslash}p{1.1cm} >{\RaggedRight\arraybackslash}p{2.5cm} >{\RaggedRight\arraybackslash}p{2.6cm} >{\RaggedRight\arraybackslash}p{2.0cm} >{\raggedleft\arraybackslash}p{1.1cm} >{\raggedleft\arraybackslash}p{1.2cm} Y}
\toprule
Etapa & Entrada & Procedimento & Critério & Responsável & Incluídos & Não retidos & Motivos \\
\midrule
Coleta e deduplicação & 12.118 & Normalização e agrupamento por DOI ou título, com proveniência preservada. & DOI normalizado; na ausência, título normalizado sem conflito de identificador. & Script Python. & 9.542 & 2.576 & Ocorrências duplicadas entre bases ou consultas. \\
Pré-triagem e resolução & 9.542 & Regras em título e resumo; auditoria de 100; corte inicial de 3.786; decisão dos 4.276 casos de dúvida. & Classe relevante ou decisão revisada \texttt{RELEVANTE}. & Scripts Python/R; avaliador único na amostra; tabela de decisão. & 3.678 & 5.864 & 4.070 dúvidas descartadas; 796 irrelevantes; 566 sem resumo; 432 complementares. \\
Alinhamento às RQs & 3.678 & Pontuação determinística de aderência às RQs e força do resumo. & Sinais de objeto, manutenção, sustentabilidade, decisão e contexto; sem exclusão forte. & Scripts R; regras aprovadas pelo pesquisador. & 137 & 3.541 & Uso secundário, descritivo, contextual ou exclusão da análise central. \\
Auditoria qualitativa & 137 & Auditoria estruturada de título, resumo, palavras-chave e campos extraídos. & Aderência ao núcleo principal e suficiência documental. & Avaliador único, sem revisão independente. & 104 & 33 & 21 secundários; 3 apenas mapeamento; 9 excluídos, incluindo 1 pendência de texto completo descartada. \\
\bottomrule
\end{tabularx}
\fonteautor
\end{table}
}

\begin{figure}[htbp]
\centering
\begin{tikzpicture}[
  node distance=6mm and 8mm,
  etapa/.style={draw, rounded corners, align=center, text width=6.0cm, minimum height=10mm, font=\scriptsize},
  final/.style={draw=black, very thick, rounded corners, align=center, text width=6.0cm, minimum height=10mm, font=\small},
  arrow/.style={-Latex, thick}
]
\node[etapa] (p1) {Busca principal \\ 12.118 ocorrências};
\node[etapa, right=of p1] (s1) {Busca de sensibilidade IA/ML \\ 6.728 ocorrências};
\node[etapa, below=of p1] (p2) {Deduplicação por DOI e título \\ 9.542 registros únicos};
\node[etapa, below=of s1] (s2) {Deduplicação interna e contra o corpus principal \\ 359 já existentes; 4.889 novos únicos};
\node[etapa, below=of p2] (p3) {Triagem, alinhamento e auditoria documental \\ 104 registros no núcleo original};
\node[etapa, below=of s2] (s3) {Classificação RQ6 e revisão dos candidatos \\ 12 novos; cinco registros existentes promovidos};
\path (p3) -- (s3) node[midway, below=15mm, final] (f) {Núcleo temático vigente \\ \textbf{121 registros}};
\draw[arrow] (p1) -- (p2);
\draw[arrow] (p2) -- (p3);
\draw[arrow] (s1) -- (s2);
\draw[arrow] (s2) -- (s3);
\draw[arrow, rounded corners=2mm] (p3.south) |- (f.west);
\draw[arrow, rounded corners=2mm] (s3.south) |- (f.east);
\end{tikzpicture}
\caption{Fluxo integrado da busca principal e da busca complementar de sensibilidade}
\label{fig:funil}
\end{figure}

A elegibilidade e a extração do núcleo temático apoiaram-se predominantemente em título, resumo, palavras-chave e campos estruturados. Entre 37 textos completos consultados em tarefas posteriores, 19 acrescentaram evidências específicas citadas individualmente; no lote de sensibilidade, sete dos 17 textos estavam disponíveis integralmente e cinco ampliaram resultados quantitativos, desenho de validação ou limitações. A extração identificou dimensões de sustentabilidade, critérios de priorização, métodos de apoio à decisão, contexto da edificação, ODS, ESG e lacunas relacionadas a instituições públicas de ensino superior, tratando cada categoria com respostas múltiplas e um único identificador por estudo. Dimensão e critério são categorias distintas, pois a dimensão é a categoria mais ampla (técnica-operacional, institucional, ambiental, de ciclo de vida, econômica ou social) e o critério é um atributo operacional específico dentro dela, de modo que energia, risco, conforto e segurança contam como critérios sem duplicar as seis dimensões. Definições e exemplos estão no manual de codificação mantido no repositório.

\subsection{Busca complementar de sensibilidade para inteligência artificial e aprendizado de máquina}
\label{sec:buscaia}

As quatro expressões booleanas do núcleo original recuperaram inteligência artificial e aprendizado de máquina apenas de forma incidental, pois nenhuma das questões originalmente formuladas, de RQ0 a RQ5, continha dicionário de termos dedicado ao tema, e o método de decisão associado à RQ3 abrangia exclusivamente instrumentos multicritério como MCDM, AHP e TOPSIS. Para verificar se essa lacuna subestimava a presença do tema no campo, uma busca complementar foi conduzida nas mesmas três bases, com expressões voltadas explicitamente a IA e aprendizado de máquina aplicados a manutenção e gestão predial, recuperando 6.728 ocorrências brutas (Tabela~\ref{tab:estrategiabusca}). Os registros foram normalizados no mesmo esquema de campos do núcleo original e deduplicados em cascata (DOI normalizado, identificador próprio da fonte, título normalizado exato e título aproximado confirmado por ano, autor ou periódico); do total bruto, 359 registros já constavam do corpus original e 4.889 registros únicos seguiram para auditoria temática.

A RQ6 recebeu manual de codificação próprio em dois níveis, com expressões que confirmam isoladamente uma técnica computacional, como \textit{machine learning}, \textit{neural network}, \textit{random forest} e \textit{large language model}, e termos ambíguos que exigem contexto adicional, como manutenção preditiva, BIM, \textit{digital twin}, IoT e \textit{decision tree} fora do contexto de algoritmo treinado. Essa distinção evita tratar manutenção preditiva, BIM ou \textit{digital twin} como sinônimos de IA, pois plataformas digitais, modelos treinados e mecanismos de decisão ocupam níveis analíticos distintos \parencite{deng_bimdigitaltwins_2021,das_iotai_2025}. Os 4.889 registros foram classificados, sem amostragem, em seis classes de evidência, com nível de confiança rebaixado por definição para registros sem resumo, concentrados no Crossref.

A classificação produziu 20 candidatos novos; a revisão individual removeu oito cuja aplicação, tecnicamente confirmada, ocorria em domínio industrial, aeroespacial ou veicular sem relação com edificações, restando 12. Sete registros já pertencentes ao corpus original, antes secundários ou excluídos por um dicionário que não reconhecia IA, foram reavaliados sob a RQ6, e cinco foram promovidos ao núcleo principal e dois permaneceram secundários por não apresentarem relação explícita com a manutenção predial. Ao final, 17 registros passaram a integrar o núcleo temático, elevando-o de 104 para 121, sob o mesmo vocabulário fechado e os mesmos procedimentos do núcleo original, sem leitura de texto completo.

\subsection{Reprodutibilidade e limites da automação}

Consultas, exportações, manual de codificação, matrizes intermediárias, scripts determinísticos, figuras e verificações automatizadas foram versionados no repositório da pesquisa, preservando identificadores, proveniência e parâmetros que permitem recompilar resultados e produtos bibliométricos, contribuição de ciência aberta que não substitui pré-registro nem revisão independente. O caráter determinístico garante que a mesma entrada, processada com a mesma versão do código, produza a mesma saída, com a subjetividade explicitada no manual de codificação, nos limiares e nas decisões sobre casos duvidosos. O desenho priorizou fontes com metadados interoperáveis, o que delimita a cobertura de relatórios técnicos, normas internas e experiências institucionais da literatura cinzenta.

Na preparação deste manuscrito para submissão, utilizou-se a ferramenta Claude (Claude Code, Anthropic, modelo Claude Sonnet 5) exclusivamente na revisão ortográfica, textual e de adequação formal às normas editoriais do periódico, sem uso na concepção, coleta, análise ou geração de conteúdo científico.
